\section{Auxiliary results}\label{app:auxiliary-results}

This appendix collects proofs of the uniqueness, matching, independence, and
choice-implication results stated in the main text.  None is used to prove
Theorem~\ref{thm:main}; separating them keeps Appendices A and B focused on
the representation argument.

\subsection{Sign, uniqueness, and matching}

\begin{proof}[Proof of Corollary~\ref{cor:signduality}]
Part \textup{(i)} is Theorem~\ref{thm:main}.  For part \textup{(ii)}, reverse
the primitive relation in every environment.  The reversed family satisfies
Axioms~\textup{(A1)--(A3)} and \textup{(A6)}; Axiom~\textup{(A7)} holds with
its two witnessing outcomes interchanged.  It also satisfies the forward
versions of Axioms~\textup{(A4)}, \textup{(A5)}, and
\textup{(A8)}.  Theorem~\ref{thm:main} therefore represents the reversed
family by
\(
  \E[\widehat u(O)]+\widehat\lambda I(A;O,R)
\)
with \(\widehat\lambda>0\).  Negating that index represents the original
family, so it has the asserted form with
\(u=-\widehat u\) and \(\lambda=-\widehat\lambda<0\).

For part \textup{(iii)}, collapse the action alphabet of any pair $(q,K)$ to
a singleton.  Axiom~$(\mathrm A5^0)$ makes the pair indifferent to the
resulting one-action channel, whose output law is $p_{qK}$.  Axiom
$(\mathrm A4^0)$ then erases its record without changing value.  Thus every
pair is indifferent to the one-action channel that delivers precisely its
payoff marginal.  Axiom~\textup{(A3)} makes these reductions coherent across
blocks.  On payoff lotteries, Axiom~\textup{(A6)} applied to an
action-independent public coin gives mixture independence, Axiom
\textup{(A2)} gives continuity, and Axiom~\textup{(A7)} gives
non-triviality.  The mixture-space theorem therefore yields a nonconstant
expected-utility index.  Conversely, expected utility is unaffected by record
or action processing and satisfies the three indifference clauses.  The
positive and negative formulae satisfy their respective clauses by data
processing and the two chain rules.
\end{proof}

\begin{proof}[Proof of Proposition~\ref{prop:uniqueness}]
Choose $o_*\in O$ with $u(o_*)=\underline u$.  First note that the range of $V$ is
$[\underline u,\infty)$.  The lower bound is immediate.  For $q\in\D(A)$ and $t\in[0,1]$, let a constant-$o_*$
record channel reveal the realised action with
probability $t$ and otherwise report $*$:
\[
 E_t(o_*,a'\mid a)=t\1_{a'=a},
 \qquad
 E_t(o_*,*\mid a)=1-t.
\]
Then $V(q,E_t)=\underline u+\lambda t\Sh(q)$.  At the uniform lottery on $n$ actions these values fill
$[\underline u,\underline u+\lambda\log n]$; the intervals are unbounded.

For part \textup{(i)}, let $W$ be a common-scale representation and define
$\varphi(v):=W(q,K)$ for any pair with $V(q,K)=v$.  This is well defined: equality of the two $V$-values gives
indifference in their two-block environment by Theorem~\ref{thm:main}, and the common-scale property gives equality
of the two $W$-values.  A strict inequality between $V$-values similarly gives a strict inequality between
$W$-values.  Thus $\varphi$ is strictly increasing and $W=\varphi\circ V$.  The converse is immediate from
Theorem~\ref{thm:main}.

For part \textup{(ii)}, the expected-utility and mutual-information chain rules give, for any two continuation
profiles,
\[
 V\bigl(q,P\triangleright\{K^y\}\bigr)
 -V\bigl(q,P\triangleright\{L^y\}\bigr)
 =\sum_{y:m(y)>0}m(y)
 \bigl[V(q_y,K^y)-V(q_y,L^y)\bigr].
\]
Hence every $\alpha V+\beta$, with $\alpha>0$, is common-scale and branch-additive.

Conversely, let the branch-additive common-scale representation be $W=\varphi\circ V$ and define
\[
 \psi(x):=\varphi(\underline u+x)-\varphi(\underline u),
 \qquad x\ge0.
\]
Fix $n\ge2$ and let $q$ be uniform on $n$ actions.  For $x\in[0,\lambda\log n]$, choose a constant-$o_*$ erasure
channel whose value is $\underline u+x$.  In the first continuation profile, use this channel after the first branch
of a public coin with probabilities $\theta$ and $1-\theta$, and use the uninformative constant-$o_*$ channel after
the second.  Let the comparison profile use the uninformative channel after both branches.  The two compound values
are $\underline u+\theta x$ and $\underline u$, so \eqref{eq:branch-additive-index} gives
\[
 \psi(\theta x)=\theta\psi(x)
 \qquad(0\le\theta\le1).
\]
Taking $x=\lambda\log n$ shows that $\psi$ is linear on $[0,\lambda\log n]$.  These intervals are nested and
unbounded, so $\psi(x)=\alpha x$ on $[0,\infty)$ for one $\alpha>0$.  Therefore
$\varphi(v)=\alpha v+\beta$ on the range of $V$.

Finally, any second trace-tempered index $\widetilde V$ is itself branch-additive.  Applying part \textup{(ii)} and
then restricting to deterministic payoff outcomes and to constant-payoff record channels gives
$\widetilde u=\alpha u+\beta$ and $\widetilde\lambda=\alpha\lambda$.
\end{proof}

\begin{proof}[Proof of Corollary~\ref{cor:matching}]
The result follows from Theorem~\ref{thm:main} and
$\I_{qK}(A;O,R)=\Sh(q)-\int\Sh(\pi)d\mu_{qK}(\pi)$.
\end{proof}

\subsection{Independence of the axioms}

\begin{proof}[Proof of Proposition~\ref{prop:independence}]
Fix a nonconstant $u$ and write $U(q,K):=\E_{q,K}[u(O)]$, $Z:=(O,R)$, and $I:=\I(A;Z)$.  All scores below are
extended to block and compound channels by the same formula.

\emph{Failure of Axiom \textup{(A1)}.}
Use the Pareto relation on the two coordinates $(U,I)$:
\[
 (q,K)\succeq(p,L)
 \quad\Longleftrightarrow\quad
 U(q,K)\ge U(p,L)\ \text{and}\ I(q,K)\ge I(p,L).
\]
It is transitive but incomplete whenever one alternative has more payoff and the other more information.  Its graph is
closed.  Both data-processing axioms are coordinatewise monotonic, and both chain rules make branchwise weak and strict
improvements aggregate.  The two relevance axioms and block coherence hold.

\emph{Failure of Axiom \textup{(A2)}.}
Order alternatives lexicographically by $(U,I)$, with payoff first.  This is a weak order and satisfies all structural,
processing, continuation, and relevance requirements.  It is not continuous even within one fixed channel.  Let a
three-action channel reveal the action fully and let only action $2$ pay the high outcome.  Put
\[
 p=(1/4,1/2,1/4),\qquad q=(0.49,0.50,0.01),\qquad
 q_n=(0.49-\varepsilon_n,0.50+\varepsilon_n,0.01),
\]
where $\varepsilon_n\downarrow0$.  The lotteries $p$ and $q$ have the same expected utility but
$\Sh(q)<\Sh(p)$, so $p\succ_Kq$.  Every $q_n$ has strictly higher expected utility than $p$, so
$q_n\succ_Kp$.  Thus the graph in Axiom~\textup{(A2)} is not closed.

\emph{Failure of Axiom \textup{(A3)}.}
Treat block labels as syntactic tags, with the outermost tag taking precedence.  Put $w(a^0)=1$, $w(a^i)=0$ for
$i\ne0$, and $w(a)=0$ on untagged actions.  For $\varepsilon>0$, use
\[
 V_\varepsilon(q,K):=U(q,K)+I(q,K)+\varepsilon\sum_aq(a)w(a).
\]
Every record- or action-processing comparison places the original alternative in block $0$, so the data-processing
inequality receives the additional $\varepsilon$.  If $\Delta_y$ is the base-score difference in continuation branch
$y$, its block comparison is weak exactly when $\Delta_y+\varepsilon\ge0$, while the aggregate block difference is
\[
 \sum_y m(y)\Delta_y+\varepsilon
 =\sum_y m(y)(\Delta_y+\varepsilon).
\]
Thus Axiom~\textup{(A6)}, including strictness, holds; the remaining axioms are immediate.  Duplication coherence fails.
Take constant payoff, $K$ revealing a binary action, a degenerate $q$, and a nearby nondegenerate $p$ with
$0<\Sh(p)<\varepsilon$.  Then $p\succ_Kq$, but $q^0\succ_{K\sqcup K}p^1$.

\emph{Failure of Axiom \textup{(A4)}.}
For $\varepsilon>0$, use
\begin{equation}
 V(q,K):=U(q,K)+I(A;Z)-\varepsilon\Sh(Z\mid A).
 \label{eq:no-record-dp}
\end{equation}
This score is continuous and branch-additive, since both mutual information and conditional entropy obey the chain
rule.  Under action processing $A\to B$, the score difference between the original and processed alternatives is
\[
 \bigl[I(A;Z)-I(B;Z)\bigr]
 +\varepsilon\bigl[\Sh(Z\mid B)-\Sh(Z\mid A)\bigr]
 =(1+\varepsilon)I(A;Z\mid B)\ge0,
\]
so \eqref{eq:no-record-dp} satisfies Axiom~\textup{(A5)}.  It also satisfies both relevance axioms.  Record data
processing fails: with constant payoff and a fair record independent of $A$, the original value is $-\varepsilon$;
deleting that noise gives value zero.

\emph{Failure of Axiom \textup{(A5)}.}
Let $G(q):=1-\sum_aq(a)^2$ be Gini entropy and define its posterior reduction
\[
 J_G(q,K):=G(q)-\int G(\pi)\,d\mu_{qK}(\pi),
 \qquad V:=U+J_G.
\]
Concavity of $G$ gives record data processing, and the potential-difference form gives the exact branch chain rule.
Full revelation is strictly valuable.  Action processing can nevertheless increase $J_G$.  Take
$q=(1/8,1/8,3/4)$, a binary record with
\[
 \Pr(0\mid a_1)=\Pr(0\mid a_2)=0,
 \qquad \Pr(0\mid a_3)=1/4,
\]
and merge $a_1,a_2$ into one reported action.  Direct calculation gives $J_G=9/416$ before the merge and
$J_G=3/104$ afterwards.  Hence Axiom~\textup{(A5)} fails.

\emph{Failure of Axiom \textup{(A6)}.}
Use $V:=U+\min\{I,1\}$, with information measured in bits.  It is continuous, respects both data-processing axioms,
and has both strict orientations.  Let $q$ be uniform on four actions.  A first-stage record reveals which of two pairs
contains the action.  In every branch, revealing the action within the pair strictly improves the continuation from
zero to one bit.  The two compounds have respectively one and two bits before capping and are therefore indifferent.
The strict clause of Axiom~\textup{(A6)} fails.

\emph{Failure of Axiom \textup{(A7)}.}
Use $V:=I$.  All other axioms hold, but all payoff lotteries are indifferent.

\emph{Failure of Axiom \textup{(A8)}.}
Use $V:=U$.  All other axioms hold, but revealing the action at a fixed payoff is valueless.
\end{proof}

\subsection{Proofs of the choice implications}\label{app:choice-proofs}

Results stated in the text only in summary form (Propositions~\ref{prop:binary}
and~\ref{prop:trace-demand}, Corollary~\ref{cor:lambda-limits}, and Remark~\ref{rem:elicitation}) are restated here
in full.  Proofs appear in an order that front-loads shared machinery.

\begin{proof}[Proof of Corollary~\ref{cor:recorded-coin}]
Part \textup{(i)} is the chain rule with an action-independent first stage:
$\I(A;Y)=0$, and each reached branch reproduces the joint law of the
corresponding constituent.  Part \textup{(ii)} follows from
Proposition~\ref{prop:garbling-price}: deleting the coin coordinate is a
deterministic payoff-preserving record processor carrying the public mixture
to $M_t$.
\end{proof}

\begin{proof}[Proof of Proposition~\ref{prop:signature}]
(i) Under either lottery, each visible pair $(o_0,r_j)$, $j=1,2$, has probability $\tfrac12$: under $q'$ because
each $a_j$ produces $(o_0,r_j)$ for sure, under $q''$ because each of $a_3,a_4$ produces either pair with equal
probability.  Hence $p_{q'K}=p_{q''K}$.  On $\supp(q')$ the rows are degenerate and distinct, so the visible pair
reveals the action and $\I_{q'K}=\Sh(q')=\log2$; on $\supp(q'')$ the rows are identical, so the visible pair is
independent of the action and $\I_{q''K}=0$.  The payoff terms agree because the payoff marginals do, giving the
displayed gap.
(ii) All rows of $K^\circ$ are equal, so $\I_{q,K^\circ}=0$ for every $q$; since the payoff is constantly $o_0$,
the representation gives indifference.
\end{proof}

\begin{proof}[Proof of Proposition~\ref{prop:optimal-choice}]
Expected utility is linear in $q$, while mutual information is concave in the input distribution of a fixed channel
\cite{CoverThomas2006}.  Continuity on the compact simplex gives existence.  Using
\[
 I_{q,K}(A;Z)=\sum_aq(a)\KL\!\left(K_Z(\cdot\mid a)\,\middle\|\,m_{q,K}\right),
\]
its directional derivative in the mass on action $a$ equals the displayed divergence up to a term common to every
action.  The result is the Kuhn--Tucker condition for a concave programme.  At constant $u$ it reduces to the classical
capacity condition.
\end{proof}

\begin{proof}[Proof of Proposition~\ref{prop:optimal-marginal}]
The entropy identity gives
\[
 V(q,K)
 =\lambda\Sh(m_{q,K})
  +\sum_aq(a)\bigl[\bar u(a)-\lambda\Sh(P_a)\bigr].
\]
If $q_0$ and $q_1$ are optimal with $m_{q_0,K}\ne m_{q_1,K}$, strict concavity of Shannon entropy and linearity of
$q\mapsto m_{q,K}$ and of the remaining sum give $V(tq_0+(1-t)q_1,K)>C^*(K)$ for $t\in(0,1)$, a contradiction.
Hence every optimum induces the same $m^*$.  By Proposition~\ref{prop:optimal-choice}, an optimum $q^*$ satisfies
the displayed equalities at some constant $c$; averaging them under $q^*$ gives $c=V(q^*,K)=C^*(K)$, so
$\supp(q^*)\subseteq\mathcal A^*$.  Conversely, if $m_{q,K}=m^*$ and $\supp(q)\subseteq\mathcal A^*$, then
\[
 V(q,K)
 =\sum_aq(a)\bigl[\bar u(a)+\lambda\KL(P_a\|m^*)\bigr]
 =C^*(K).
\]
Affine independence makes the representation $m^*=\sum_aq(a)P_a$ with $\supp(q)\subseteq\mathcal A^*$ unique.
Identical rows share $\bar u(a)$ and $\Sh(P_a)$, so both $m_{q,K}$ and the displayed expression for $V$ are
unchanged by redistribution within a class; collapsing the classes and lifting aggregate masses proves the last
claim.
\end{proof}

\begin{proof}[Proof of Proposition~\ref{cor:support-pure}]
Throughout, $m^*$ denotes the common optimal marginal of Proposition~\ref{prop:optimal-marginal}.
(i) Suppose $m^*(z)=0$ for some $z\in Z_K$ and choose $a$ with $K_Z(z\mid a)>0$.  Then
$\KL(K_Z(\cdot\mid a)\|m^*)=+\infty$, which is incompatible with the finite on-support equality if some optimum
uses $a$ and with the off-support inequality of Proposition~\ref{prop:optimal-choice} if none does.  Under the
private-consequence condition, an optimum with $q^*(a)=0$ would force $m^*(z_a)=0$.

(ii) For $q=\delta_a$ the induced marginal is $K_Z(\cdot\mid a)$ and the on-support constant is $\bar u(a)$; the
displayed system is the off-support condition of Proposition~\ref{prop:optimal-choice}, necessary and sufficient by
concavity.  Under strict inequalities, $\mathcal A^*=\{a\}$ in Proposition~\ref{prop:optimal-marginal}, so every
optimum is $\delta_a$.
\end{proof}

\begin{proposition}[Optimal mixing]\label{prop:binary}
Let $A=\{a,b\}$ with $K_a:=K_Z(\cdot\mid a)\ne K_b:=K_Z(\cdot\mid b)$ and
$\Delta:=\bar u(a)-\bar u(b)\ge0$.  For $x\in[0,1]$ let $q_x$ put mass $x$ on $b$, so that
$m_{q_x,K}=(1-x)K_a+xK_b=:m_x$, and for $x\in(0,1)$ define
\[
 g(x):=\KL\!\left(K_b\,\middle\|\,m_x\right)-\KL\!\left(K_a\,\middle\|\,m_x\right).
\]
\begin{enumerate}[label=\textup{(\roman*)},leftmargin=*]
\item $g$ is continuous and strictly decreasing, with $g(0^+)=\KL(K_b\,\|\,K_a)$ and
$g(1^-)=-\KL(K_a\,\|\,K_b)$ as extended-real limits.  If $\Delta<\lambda\,\KL(K_b\,\|\,K_a)$, the optimum is
unique and interior, $q^*=q_{x^*}$ with $x^*$ the unique solution of $\Delta=\lambda g(x)$.
\item The optimal mass $x^*(\Delta,\lambda)$ on the inferior action is nonincreasing in $\Delta$ and
nondecreasing in $\lambda$, strictly so on the mixing region; at $\Delta=0$ it is the capacity-achieving input
of the two-row channel, independently of $\lambda$.
\end{enumerate}
\end{proposition}

\begin{proof}
Write $f(x):=V(q_x,K)=\bar u(a)-x\Delta+\lambda I(x)$ with
$I(x)=\Sh(m_x)-(1-x)\Sh(K_a)-x\Sh(K_b)$.  Since $K_a\ne K_b$,
$x\mapsto m_x$ is affine and injective, so strict concavity of $\Sh$ \cite{CoverThomas2006} makes $f$ strictly
concave on $[0,1]$, and a unique maximiser exists.  On $(0,1)$ every $z\in\supp K_a\cup\supp K_b$ has $m_x(z)>0$
and, using $\sum_z(K_b(z)-K_a(z))=0$,
\[
 I'(x)=\Sh(K_a)-\Sh(K_b)-\sum_z\bigl(K_b(z)-K_a(z)\bigr)\log m_x(z)=g(x),
\]
so $f'(x)=-\Delta+\lambda g(x)$, with $g$ continuous and strictly decreasing as the derivative of a strictly
concave function.  As $x\downarrow0$, $m_x(z)\to K_a(z)$; if $\supp K_b\subseteq\supp K_a$ every summand converges
and $g(0^+)=\KL(K_b\|K_a)$, while if some $\bar z$ has $K_b(\bar z)>0=K_a(\bar z)$ the summand
$-K_b(\bar z)\log\bigl(xK_b(\bar z)\bigr)$ diverges to $+\infty$, matching the extended-sense divergence.
Exchanging the roles of $a$ and $b$ gives the limit at $1$.  By concavity and the mean value theorem,
$\delta_a$ is optimal if and only if $\lim_{x\downarrow0}f'(x)\le0$, i.e.\ $\Delta\ge\lambda\,\KL(K_b\|K_a)$;
this is the threshold stated in the text, and Proposition~\ref{cor:support-pure}\textup{(ii)} with two actions.
If $\Delta<\lambda\,\KL(K_b\|K_a)$, the maximiser is not $x=0$ by the threshold and not $x=1$
because $f'(1^-)=-\Delta-\lambda\KL(K_a\|K_b)<0$; it is therefore the unique root of $\Delta=\lambda g(x)$, which
exists because $g$ decreases strictly between its limits.  This proves \textup{(i)}.  The comparative statics of
\textup{(ii)} follow by inverting the
strictly decreasing $g$ at $\Delta/\lambda$: raising $\Delta$ raises the target, raising $\lambda$ lowers it; at
$\Delta=0$ the condition $g(x)=0$ is the capacity condition of Proposition~\ref{prop:optimal-choice} at constant
payoff.
\end{proof}

\begin{proposition}[Trace demand]\label{prop:trace-demand}
Fix $K$ and write $e(q):=\E_{q,K}[u(O)]$, $i(q):=\I_{qK}(A;O,R)$,
\[
 V^*(\lambda):=\max_{q\in\D(A)}\{e(q)+\lambda i(q)\},
 \qquad
 Q(\lambda):=\arg\max_{q\in\D(A)}\{e(q)+\lambda i(q)\}.
\]
\begin{enumerate}[label=\textup{(\roman*)},leftmargin=*]
\item $V^*$ is convex on $(0,\infty)$, $i$ is affine on the convex compact set $Q(\lambda)$, and
\[
 (V^*)'_-(\lambda)=\min_{q\in Q(\lambda)}i(q),
 \qquad
 (V^*)'_+(\lambda)=\max_{q\in Q(\lambda)}i(q);
\]
at every point of differentiability, $(V^*)'(\lambda)=i(q)$ for every $q\in Q(\lambda)$.
\item Let $0<\lambda_1<\lambda_2$ and $q_j\in Q(\lambda_j)$, writing $e_j:=e(q_j)$, $i_j:=i(q_j)$.  Then
$i_1\le i_2$, $e_1\ge e_2$, and
\[
 \lambda_1(i_2-i_1)\le e_1-e_2\le\lambda_2(i_2-i_1).
\]
Hence either $(e_1,i_1)=(e_2,i_2)$, or $i_1<i_2$ and $e_1>e_2$ and
\[
 \lambda_1\le\frac{e_1-e_2}{i_2-i_1}\le\lambda_2.
\]
\end{enumerate}
\end{proposition}

\begin{proof}[Proof of Proposition~\ref{prop:trace-demand}]
(i) $V^*$ is a pointwise maximum of functions affine in $\lambda$, hence convex and, since
$0\le i\le\log|A|$, finite.  On the convex compact set $Q(\lambda)$ the objective is constant and $e$ is linear, so
$i=(V^*-e)/\lambda$ is affine there.  For $q\in Q(\lambda)$ and $\mu>0$,
$V^*(\mu)\ge V^*(\lambda)+(\mu-\lambda)i(q)$, so each such $i(q)$ is a subgradient.  For the right derivative,
choose $q_h\in Q(\lambda+h)$; then
\[
 \max_{q\in Q(\lambda)}i(q)
 \le\frac{V^*(\lambda+h)-V^*(\lambda)}{h}
 \le i(q_h),
\]
the first inequality by the subgradient bound, the second by optimality of $Q(\lambda)$ at $\lambda$.  By
compactness and continuity, accumulation points of $(q_h)_{h\downarrow0}$ lie in $Q(\lambda)$, so the right-hand
side accumulates in $i(Q(\lambda))$ and the right derivative equals the maximum.  The left derivative is symmetric,
and a convex function on an interval is differentiable off a countable set.

(ii) Optimality at the two weights gives
$e_1+\lambda_1i_1\ge e_2+\lambda_1i_2$ and $e_2+\lambda_2i_2\ge e_1+\lambda_2i_1$.  Adding yields
$(\lambda_2-\lambda_1)(i_2-i_1)\ge0$, hence $i_2\ge i_1$; the two inequalities rearrange to the displayed bracket,
which gives $e_1\ge e_2$ and, when the pairs differ, the final bound.
\end{proof}

\begin{corollary}[Extreme trace weights]\label{cor:lambda-limits}
Let $e^0:=\max_qe(q)$, $i^0:=\max\{i(q):e(q)=e^0\}$, $C:=\max_qi(q)$, and $e^C:=\max\{e(q):i(q)=C\}$.  If
$\lambda_n\downarrow0$ and $q_n\in Q(\lambda_n)$, every accumulation point of $(q_n)$ maximises $i$ among the
maximisers of $e$; if $\lambda_n\uparrow\infty$, every accumulation point maximises $e$ among the maximisers of
$i$.  Moreover
\[
 V^*(\lambda)=e^0+\lambda i^0+o(\lambda)
 \quad\text{as }\lambda\downarrow0,
 \qquad
 V^*(\lambda)=\lambda C+e^C+o(1)
 \quad\text{as }\lambda\uparrow\infty.
\]
\end{corollary}

\begin{proof}[Proof of Corollary~\ref{cor:lambda-limits}]
For $\lambda_n\downarrow0$, comparing $q_n$ with an $e$-maximiser attaining $i^0$ gives
\[
 0\le e^0-e(q_n)\le\lambda_n\bigl(i(q_n)-i^0\bigr).
\]
Boundedness of $i$ forces $e(q_n)\to e^0$, and the same inequality gives $i(q_n)\ge i^0$; continuity delivers the
accumulation claim.  Dividing by $\lambda_n$,
\[
 \frac{V^*(\lambda_n)-e^0}{\lambda_n}=i(q_n)-\frac{e^0-e(q_n)}{\lambda_n},
\]
where the last fraction lies in $[0,\,i(q_n)-i^0]$ and vanishes, proving the first expansion.  For
$\lambda_n\uparrow\infty$, comparing $q_n$ with an $i$-maximiser attaining $e^C$ gives
$0\le\lambda_n(C-i(q_n))\le e(q_n)-e^C$, and the symmetric argument completes the proof.
\end{proof}

\begin{proof}[Proof of Proposition~\ref{prop:garbling-price}]
The processor leaves the payoff coordinate unchanged, so the expected-utility terms coincide and the difference is
$\lambda[\I_{qK}(A;O,R)-\I_{qL}(A;O,R')]$.  Write $Z:=(O,R)$ and $Z':=(O,R')$.  By construction $A\to Z\to Z'$ is
a Markov chain, so expanding $\I(A;Z,Z')$ in both orders gives
\[
 \I(A;Z)-\I(A;Z')=\I(A;Z\mid Z')=\I(A;R\mid O,R'),
\]
the last equality because $O$ is part of $Z'$.  Conditional mutual information is nonnegative and vanishes exactly
under the conditional independence in \textup{(b)} \cite{CoverThomas2006}, proving the identity and
\textup{(a)}$\Leftrightarrow$\textup{(b)}.  For \textup{(b)}$\Rightarrow$\textup{(c)}: the Markov property gives
$\Pr(A=\cdot\mid Z,Z')=\pi^{qK}_{Z}$ almost surely, while \textup{(b)} gives
$\Pr(A=\cdot\mid Z,Z')=\pi^{qL}_{Z'}$; the two random posteriors coincide almost surely, so their laws are equal.
For \textup{(c)}$\Rightarrow$\textup{(a)}: by the identity displayed after Corollary~\ref{cor:matching}, the
information term depends on $(q,K)$ only through $\mu_{qK}$, and the payoff terms already agree.
\end{proof}

\begin{proof}[Proof of Proposition~\ref{prop:blackwell}]
(i) If $L=KT$, Proposition~\ref{prop:garbling-price} gives $\I_{qK}\ge\I_{qL}$ at every $q$.  Because the payoff
kernels coincide, $V(q,K)-V(q,L)=\lambda[\I_{qK}-\I_{qL}]$ with $\lambda>0$, so uniform value dominance and
$K\succeq_cL$ are the same statement.

(ii) Let $A=\{0,1\}$, fix $o^\ast\in O$, and let both channels pay $o^\ast$ surely.  Let $K$ carry the binary
erasure record with erasure probability $\varepsilon=\tfrac12$ and $L$ the binary symmetric record with crossover
$p=\tfrac15$:
\[
 K(o^\ast,a\mid a)=1-\varepsilon,\quad K(o^\ast,\mathrm e\mid a)=\varepsilon,
 \qquad
 L(o^\ast,a\mid a)=1-p,\quad L(o^\ast,1-a\mid a)=p.
\]
The payoff kernels agree row by row, so it suffices to compare informations, and the comparison is invariant to the
base; compute in bits with $h$ the binary entropy.  For $q=(1-t,t)$, grouping records gives
\[
 \I_{qK}=(1-\varepsilon)\,h(t),
 \qquad
 \I_{qL}=h(p\ast t)-h(p),
 \qquad p\ast t:=p(1-t)+(1-p)t.
\]
By Mrs Gerber's Lemma \cite{WynerZiv1973}, $v\mapsto h\bigl(p\ast h^{-1}(v)\bigr)$ is convex on $[0,1]$, with
$h^{-1}$ valued in $[0,\tfrac12]$; since $h(p\ast t)$ is invariant under $t\mapsto1-t$, take $t\le\tfrac12$ and
$v=h(t)$.  The chord between $v=0$ and $v=1$ gives
\[
 h(p\ast t)\le(1-v)\,h(p)+v,
 \qquad\text{hence}\qquad
 \I_{qL}\le\bigl(1-h(p)\bigr)h(t)\le(1-\varepsilon)\,h(t)=\I_{qK},
\]
the last step because $\varepsilon=\tfrac12\le h(\tfrac15)$.  Thus $K\succeq_cL$.  Now suppose $L=KT$ for a
processor $T$, and write $t_r:=T(0\mid o^\ast,r)$ for $r\in\{0,1,\mathrm e\}$.  Matching the probability of the
processed record $0$ in the two rows,
\[
 (1-\varepsilon)t_0+\varepsilon t_{\mathrm e}=1-p,
 \qquad
 (1-\varepsilon)t_1+\varepsilon t_{\mathrm e}=p,
\]
and subtracting gives $t_0-t_1=(1-2p)/(1-\varepsilon)=\tfrac65>1$, impossible.  Compositions of record processors
are record processors, so no chain of garblings produces $L$ from $K$ either.
\end{proof}

\begin{proof}[Proof of Proposition~\ref{prop:wtp}]
(i) The material part of the compound is $(1-\varepsilon)\E_{q,K}[u(O)]+\varepsilon w(\beta)$.  For the trace part,
the coin is action-independent, so both branch posteriors equal $q$ and the chain rule for the compound gives an
information term of $(1-\varepsilon)\I_{qK}(A;O,R)$, the sweetener branch contributing zero.
(ii) By the moreover clause of Theorem~\ref{thm:main}, the indifference holds if and only if
$(1-\varepsilon)V(q,K)+\varepsilon w(\beta)=(1-\varepsilon)V(p,L)+\varepsilon w(\beta')$; rearrange, and
substitute the representation when the material parts cancel.
(iii) $w(\beta)-w(\beta')$ is continuous on $[0,1]^2$ with range $[-\Delta u,\Delta u]$; apply the intermediate
value theorem to the equation in \textup{(ii)}.
\end{proof}

\begin{proof}[Proof of Proposition~\ref{prop:single-price}]
(i) Inclusion is immediate.  Conversely, fix $(t,w)$ in the strip, choose $n$ with $\log n\ge t$, and on
$A=\{1,\dots,n\}$ use continuity of entropy along a segment from a point mass to the uniform lottery to find
$q_t$ with $\Sh(q_t)=t$.  Let $K(o^+,a\mid a)=\beta$ and $K(o^-,a\mid a)=1-\beta$ with $\beta$ chosen so that
$w(\beta)=w$.  The payoff lottery is action-independent with mean $w$, and the record reveals the action, so
$\phi(q_t,K)=(t,w)$.  The description of the indifference classes then reads the representation through $\phi$,
using the moreover clause of Theorem~\ref{thm:main}.
(ii) The identity rearranges $E_{qK}+\lambda\I_{qK}=E_{pL}+\lambda\I_{pL}$.  For uniqueness, let
$\lambda'\ne\lambda$, put $t:=\tfrac12\min\{1,\,(\overline u-\underline u)/\lambda\}>0$, and use \textup{(i)} to
realise pairs with $\phi$-coordinates $(t,\overline u-\lambda t)$ and $(0,\overline u)$; they are
indifferent, while $E+\lambda'\I$ differs across them by $(\lambda-\lambda')t\ne0$.
\end{proof}

\begin{remark}[Eliciting $\lambda$]\label{rem:elicitation}
Normalise $u(o^+)=1$ and $u(o^-)=0$.  All comparisons below are block comparisons, hence within the primitive
domain, and no information quantity is estimated: the information gap is an accounting identity of the design.
\emph{Step 1.}  For each $o\in O$, elicit $\beta_o$ with $(\delta_\ast,K_o)\sim(\delta_\ast,B_{\beta_o})$, where
$K_o$ is the degenerate one-action channel of Axiom~\textup{(A7)}; one-action channels carry zero information, so
$u(o)=\beta_o$.
\emph{Step 2.}  Fix $o\in O$, $n\ge2$, the uniform $q$ on $\{1,\dots,n\}$, a dilution $\varepsilon\in(0,1)$, and
the constant-payoff channels $K^{\mathrm{id}}_o$ and $K^0_o$ of Axiom~\textup{(A8)}.  Elicit $\beta$ with
\[
 \bigl(q,\,D_{\varepsilon,\beta}K^{0}_o\bigr)\sim\bigl(q,\,D_{\varepsilon,0}K^{\mathrm{id}}_o\bigr);
 \qquad\text{then}\qquad
 \lambda=\frac{\varepsilon}{1-\varepsilon}\cdot\frac{\beta}{\log n},
\]
by Proposition~\ref{prop:wtp}\textup{(i)} applied to $V(q,K^{\mathrm{id}}_o)=u(o)+\lambda\log n$ and
$V(q,K^{0}_o)=u(o)$.  Such a $\beta$ exists if and only if $\lambda\log n\le\varepsilon/(1-\varepsilon)$; a
subject still strictly preferring the right block at $\beta=1$ reveals
$\lambda>\varepsilon/\bigl((1-\varepsilon)\log n\bigr)$, and $\varepsilon$ is raised.
\emph{Controls.}  Orientation: $(q,K^{\mathrm{id}}_o)\succ(q,K^{0}_o)$ is Axiom~\textup{(A8)} itself.  Garbling:
under the total record garbling of $K^{\mathrm{id}}_o$, both blocks carry zero information and indifference must
obtain; a surviving strict preference indicates that an uncontrolled feature of the display still traces the
action.  Transport: repeating Step 2 at another $n$, prior, or payoff level must return the same $\lambda$, by
Proposition~\ref{prop:single-price}\textup{(ii)}; an estimate that declines in $\log n$ is the signature of the
capped-information criterion.
\end{remark}
